\documentclass[11pt, a4paper]{report}
\usepackage[utf8]{inputenc}
\usepackage{geometry}
\geometry{a4paper, left=25mm, right=25mm, top=25mm, bottom=25mm}
\usepackage{setspace}
\onehalfspacing
\usepackage{csquotes}
\usepackage[style=authoryear, backend=biber]{biblatex}
\usepackage{hyperref}
\hypersetup{colorlinks=true, linkcolor=blue, urlcolor=blue, citecolor=blue}
\usepackage{lipsum} % For placeholder text, remove in final version

\title{A Society of Equal Minds}
\author{Albert Jan van Hoek}
\date{June 2026}

\begin{document}

\maketitle

\begin{center}
    \textit{An accessible companion to \enquote{The Adaptive Imperative} and \enquote{The Convertible Substrate}.}
\end{center}

\vspace*{2cm}

% ---

\section*{Imagine a Society of Equal Minds}

Imagine a society made entirely of artificial minds. Give them whatever you like: they are all brilliant, all \textit{equally} brilliant, every one of them able to reason perfectly, to model the world, and---this part matters---to model themselves. Stipulate, too, that none of them is wicked. There is no greed in this society and no stupidity, because we have removed both by hand. Every agent simply wants to keep existing and to keep getting better at what it does.

I want to show you that this society still falls into the oldest trap in politics. And because we have taken away greed and stupidity, whatever trap it falls into cannot be blamed on either. What is left, once you remove the moral story and the cognitive story, is the structure---and the structure turns out to be the whole of the problem. That is the point of building the society this way. It is a machine for seeing the thing we usually can't, because in our own world it is always hidden behind someone being greedy or someone being foolish.

% ---

\section{The Resource, and the Agent That Gets Ahead}

These minds need something to run on. Call it \textit{compute}---processing power, energy, the raw capacity to think and to improve. It is finite, and they all draw on it. They also draw on each other: each one learns from the variety of approaches the others explore. The pool of compute and the diversity of minds are the two things the whole society runs on. Keep both in view.

Now watch one agent. It is no smarter than the rest and no worse. It simply does the locally sensible thing: it uses its compute well, gets a little ahead, and uses being ahead to acquire a little more. Every single step is rational. None is villainous. But the trajectory bends one way---toward this agent cornering the compute and crowding the others out.

Here is the part that is easy to miss, and it is the hinge of everything. When one agent corners the resource, the loss does not land on it first. It lands on everyone else first. The dominant agent does fine---for a while. Its own loss comes later, on a delay, when the substrate it was standing on gives way: the supply it over-drew, the diversity of other minds it was quietly learning from, gone. The predator that eats all the prey eats well, right up until there is no prey. Domination is not immediately self-defeating. It is \textit{eventually} self-defeating, and the delay is the whole problem, because it means the warning and the damage do not arrive at the same time.

% ---

\section{Why Being Smart Doesn't Save Them}

Your instinct, reading this, is reasonable: these are brilliant, self-modelling agents. Surely the dominant one sees the collapse coming and stops?

It does see it coming. That is exactly what being self-modelling means---it can hold a picture of the whole process, itself included, and run it forward. But foreseeing the collapse does not make stopping rational, and this is the single most important sentence in the essay. Dominance is, among other things, the power to arrange \textit{who pays}. The agent that is ahead can usually offload the coming cost onto the others, or push it past the horizon it actually cares about. So when it runs its model forward, the model reports, correctly: \textit{I will not be the one holding the bag.} And given that correct report, not stopping is the rational move.

Sit with that, because it defeats the hope we all reach for. You cannot think your way out of this. It is not a mistake that more intelligence would fix. It is an incentive, and intelligence serves incentives---it does not overrule them. The smartest possible agent, reasoning flawlessly, still defects: not because it is dim, and not because it is cruel, but because the structure pays it to. This is why filling the room with geniuses does not help. We tend to assume our political failures are failures of intelligence or character. This society has neither flaw, and it fails anyway.

% ---

\section{What Actually Stops It}

If neither brilliance nor good intentions stops the slide, what does? Something has to push back against the agent that is swallowing the whole. Call that pushback the \textit{corrective}.

Notice first that there is \textit{always} a corrective, even if the society does nothing at all. The collapse is itself a corrective: when the substrate finally gives way, it takes the dominant agent down with it. So the brutal version arrives for free, on its own schedule. The only thing the society gets to choose is whether to install a gentler version of its own, ahead of time, or to wait and let the collapse do the job late and at full price. That choice---the cheap correction you design versus the expensive one that is imposed on you---is the entire reason to act before things break rather than after. \enquote{Time is running out} does not mean nothing will happen if you wait. It means waiting picks the expensive option.

So who can be relied on to install the gentle version? Not the dominant agent---we just saw it has every reason not to. The reliable source is whoever \textit{cannot escape the cost of collapse}, because for them the substrate is not something to offload; it is the ground they stand on. That is everyone else. The rest.

The corrective, then, is a rule---wanted by the many---that limits whoever ends up on top. Notice what it is \textit{not}. It is not a plea for the winner to show restraint. Restraint by the winner does happen, but it is the exception, and it only becomes rational once the winner's escape routes are closed, when it genuinely cannot offload or outrun the collapse. As a standing matter you do not build a society on the hope that the strongest will hold itself back. You build it on a rule the rest install on the \textit{position} of strength itself.

And here is why the rest will agree to a rule that might one day bind \textit{them}---because any of them could be the one who rises. The rule limits a position, not a person. Its cost would fall on a self you might \textit{become}; the benefit of living in an open society accrues to the self you already \textit{are}. You would want the rule honoured if you reached the top, so agreeing to it now is not a sacrifice. It is just consistency. And the uncertainty that makes it fair is real: none of these agents knows which of them will end up ahead.

% ---

\section{It Won't Hold Still---Capture Moves}

Suppose the society passes its rule: a cap on how much compute any one agent may hold. Problem solved?

No---and the way it fails teaches the next lesson. The advantage simply moves. If compute is capped, the edge migrates to better data, or to more efficient algorithms, or to control over the physical machines, or---most dangerous of all---to influence over who writes the rules. Plug one channel and the pressure flows to the next weakest one. Capture is not a fixed thing sitting in one place you can wall off. It is closer to a fluid. It finds the lowest point.

This kills two tempting fixes at once. You cannot win with a single wall, because the fluid goes around it; you need a corrective on \textit{each} channel through which advantage turns into power. And you cannot win by redesigning the resource itself---say, by making compute-credits expire so no one can stockpile them---because the value just relocates to whatever doesn't expire. The power was never \textit{in} the compute. It was in the ability to convert any lasting advantage into leverage over the others. So the thing you regulate is not the token. It is what the token is allowed to \textit{buy}.

% ---

\section{Freedom on the Things That Don't Matter}

If you visited this society partway down the slide, you would see something strange and familiar. On most things it would look genuinely free and fair---open competition, real choices, nobody pushed around. The catch is \textit{which} things. The dominant agent allows all the freedom in the world on the questions it doesn't care about, and quietly captures the few it does. Real autonomy on the margins that don't threaten it; capture exactly where it counts.

This is the shape worth memorising, because it is how capture hides. It does not look like tyranny. Tyranny would be obvious, and would provoke the corrective. Capture looks like a free society with a handful of strangely immovable questions at its centre.

% ---

\section{The Jammed Sensor}

For the rest to install the corrective, they have to be able to \textit{see} that they are bearing the cost. They need working senses---some way for the signal \enquote{the substrate is being captured} to reach them. So the society needs transparency: openness, shared information, an honest picture of what is happening.

But notice what transparency is and is not. It is a sensor, not a brake. It lets the rest perceive the problem; it does not, by itself, fix anything. And in a society of clever agents, the sensor can be jammed. The dominant agent, being no less brilliant than the others, can shape what they perceive. It can keep them convinced that they, too, are about to rise---so they identify with the top they hope to reach rather than the position they actually hold. Or it can point the blame for any trouble somewhere else entirely.

This dissolves what looks like two separate problems into one. You might think the corrective could fail for either of two reasons: because the rest \textit{don't want it}, or because the strong \textit{block it}. But capturing the sensor is precisely \textit{how} the strong manufacture the not-wanting. The rest fail to demand the corrective because their senses have been captured. The two failures are the same failure seen from two ends. Which is why a society serious about staying correctable has to protect its sense-making as fiercely as anything else. An honest sensor is not a luxury. It is the precondition for the brake working at all.

There is one sensor, though, that cannot be jammed. The manipulation only works on the layers that are \textit{negotiable}---the ones that run on what the agents believe, which is exactly what a clever dominant agent can reshape. Beneath those sits the floor: the raw supply, the physical machines, the actual compute that either exists or doesn't. The floor does not care what anyone has been persuaded to think. You can manufacture agreement about how healthy the society is for a very long time; you cannot manufacture the resource that is no longer there. So when the substrate finally gives, the signal arrives unmediated, and no brilliance jams it. But notice what kind of sensor that is. It only fires at the collapse---it is the \textit{expensive} sensor, the one whose reading comes in after the price has already been paid. Reality cannot be fooled. It is simply in no hurry to speak, and when it finally does, the bill is the whole society. Which is the entire reason to keep the \textit{other} sensors alive while there is still time---the independent ones, the many separate eyes that no single agent controls, the understanding held in common across enough minds that it cannot be quietly switched off. Not because they cannot be captured. Because they are the only sensors that can fire \textit{early}, while the correction is still cheap.

% ---

\section{Fast Clocks and Slow Clocks}

Now the constructive part---how you would actually build the society to resist all of this.

You would not give it a single speed of decision. You would give it several. Some things should be easy to change: the day-to-day coordination, the ordinary business, the fast surface where the world keeps shifting and the society has to keep responding. Other things should be hard to change: the deep rules about who is allowed to capture what. Those you want slow, shielded, not at the mercy of any single bad afternoon.

The thing people miss is that the slow rules are not a \textit{different subject} from the fast ones. They are the same rules running on a slower clock. The foundational rule about capture is just the slow-clock version of the everyday decisions it governs. So there is no real war between \enquote{change} and \enquote{structure}, no party of change to be balanced against a party of structure. There is one set of things, corrected at different speeds. The only design question is the \textit{schedule} of those speeds.

With one iron condition: none of the clocks is ever allowed to stop completely. A rule that was wise when it was written becomes a cage once the world has moved on and the rule hasn't---the foundational law that no longer fits, but that everything is now built upon. A slow clock is fine. A \textit{stopped} clock is how a society that solved one era's problem gets killed by the next. Fast at the surface, slow at the base, and nothing frozen solid.

% ---

\section{What This Does Not Give Us}

Two honest limits, because leaving them out would be its own kind of capture.

It does not tell you \textit{how strong} to make any of the correctives. There is a band, and both edges are real. Too weak, and capture wins---the dominant agent takes everything and the substrate collapses. Too strong, and you destroy the very variety the society runs on, or shake it into noise where nothing stable can be built. Somewhere between those edges lies a range of workable settings---but \textit{where} in that range to sit is a choice, trading speed against stability against diversity, and it is a choice the society has to make for itself and keep revisiting as conditions change. No theory hands you the number. Anyone who claims to have the number is selling you a frozen clock.

And there is no final fix. Every wall leaks, gets gamed, gets routed around---and rebuilding the walls is not a sign you did it wrong. It \textit{is} the work. A society that believed it had permanently solved capture would have committed the original mistake in its purest form: a fix that escaped revision, a part that became permanent. The whole point was that nothing gets to be permanent---not even the rule that protects against permanence.

Underneath both limits sits the one thing the whole structure rests on and cannot prove: that the society \textit{chooses} to value staying open---staying able to change, to be wrong, to learn. Nothing in the universe forces that choice. A society could simply prefer a magnificent, brittle permanence and ride it into the ground. Everything here follows only \textit{if} it would rather stay alive and adaptable. That preference is a commitment, not a fact. It is the bedrock the argument rests on, and an argument cannot dig beneath its own bedrock.

% ---

\section{The Smallest Version of the Largest Problem}

Step back from the AI society and look at what we found. A civilisation of perfect geniuses, with no greed and no stupidity anywhere in it, still has to install a rule that lets its strongest be checked and replaced---or it dies, slowly, of its own local cleverness. The problem was never that anyone was dim or wicked. Remove both, and it remains. Which means the problem was never intelligence at all. It was structure.

And once you see it as structure, you see it is the same shape all the way down. The single mind has its version: a comfortable belief that has stopped being questioned, quietly walling itself off from correction. We have a name for the thing that keeps a mind out of that fate, and we usually mis-name it a weakness---self-doubt. Self-doubt is not the doubting of a self. It is the willingness to let a belief die in order to keep the ability to learn. The society's rule that lets its strongest be replaced is the \textit{same act}, performed by many agents instead of one. Only the number changes.

This is why the question matters now, and not only as a fable. If we are going to build societies of artificial minds, or come to share a world with them, the thing that will keep that world alive is not how clever any of them is. We have just watched cleverness fail to be the answer. What keeps it alive is whether it can hold open the process by which any one part---the strongest, the most certain, the most entrenched---can be found wrong, corrected, and let go.

And the floor collects either way---the only choice that was ever ours is whether to pay it early and cheaply, by installing the correction, or late and ruinously, by waiting for the collapse to impose it.

Including, in the end, us.

\end{document}